% texte.tex — f169r, vv. 551–589 (conversion initiale des travaux antérieurs)
\vers[n=551]{ΜΙ.}{Mika}{ἀκούετ᾽ ὦ γυναῖκες οἷ᾽ εἴρηκεν ἡ πανοῦργος}{Vous entendez, femmes, comment la coquine}
\vers[n=552]{}{}{ἡμᾶς ἁπάσας αὖθις αὖ.}{nous arrange toutes, une fois de plus.}
\vers[suite]{ΜΝ.}{Mnésiloque}{καὶ νὴ Δί᾽ οὐδέπω γε}{Et, par Zeus, je n'ai pas encore}
\vers[n=553]{}{}{εἴρηχ᾽ ὅσα ξύνοιδ᾽· ἐπεὶ βούλεσθε πλείον᾽ εἴπω;}{dit tout ce que je sais. En voulez-vous davantage ?}
\vers[n=554]{ΜΙ.}{Mika}{ἀλλ᾽ οὐκ ἂν ἔτ᾽ ἔχοις· ὅσα γὰρ ᾔδησθ᾽ ἐξέχεας ἅπαντα.}{Tu ne pourrais. Ce que tu savais, tu l'as épanché.}
\vers[n=555]{ΜΝ.}{Mnésiloque}{μὰ Δί᾽ οὐδὲ τήν γε μυριοστὴν μοῖραν ὧν ποιοῦμεν.}{Non, par Zeus, pas le dix-millième de ce que nous faisons.}
\vers[n=556]{}{}{ἐπεὶ τάδ᾽ οὐκ εἴρηχ᾽, ὁρᾷς, ὡς στλεγγίδας λαβοῦσαι}{Car enfin, je n'ai pas dit – tu vois ? – comment nous prenons des strigiles}
\vers[n=557]{}{}{ἔπειτα σιφωνίζομεν τὸν κύσθον.}{pour ensuite nous siphonner le con...}
\vers[suite]{ΜΙ.}{Mika}{ἐπιτριβείης.}{La peste soit de toi.}
\vers[n=558]{ΜΝ.}{Mnésiloque}{ὥς τ᾽ αὖ τὰ κρέ᾽ ἐξ Ἀπατουρίων ταῖς μαστροποῖς διδοῦσαι}{... que les viandes des Apaturies, nous les distribuons à nos procureuses,}
\vers[n=559]{}{}{ἔπειτα τὴν γαλῆν φαμεν–}{puis que nous accusons la belette...}
\vers[suite]{ΜΙ.}{Mika}{τάλαιν᾽ ἐγώ· φλυαρεῖς.}{Merci de moi ! Tu radotes.}
\vers[n=560]{ΜΝ.}{Mnésiloque}{οὐδ᾽ ὡς τὸν ἄνδρα τῷ πελέκει γυνὴ κατεσπόδησεν,}{... ni que telle autre occit son époux d'un coup de hache,}
\vers[n=561]{}{}{οὐκ εἶπον· οὐδ᾽ ὡς φαρμάκοις ἑτέρα τὸν ἄνδρ᾽ ἔμηνεν,}{je ne l'ai pas dit ; ni qu'une autre rendit fou son mari par des drogues,}
\vers[n=562]{}{}{οὐδ᾽ ὡς ὑπὸ τῇ πυέλῳ κατώρυξέν ποτε—}{ni que naguère sous la baignoire...}
\vers[suite]{ΜΙ.}{Mika}{ἐξόλοιο.}{Crève...}
\vers[n=563]{ΜΝ.}{Mnésiloque}{Ἀχαρνικὴ τὸν πατέρα.}{... cette Acharnienne enfouit son père.}
\vers[suite]{ΜΙ.}{Mika}{ταυτὶ δῆτ᾽ ἀνέκτ᾽ ἀκούειν;}{Supporterons-nous ce langage ?}
\vers[n=564]{ΜΝ.}{Mnésiloque}{οὐδ᾽ ὡς σὺ τῆς δούλης τεκούσης ἄρρεν εἶτα σαυτῇ}{Ni que toi, quand ta servante mit au jour un garçon,}
\vers[n=565]{}{}{τοῦθ᾽ ὑπεβάλου, τὸ σὸν δὲ θυγάτριον παρῆκας αὐτῇ.}{tu substituas l'enfant, et livras ta petite fille.}
\vers[n=566]{ΜΙ.}{Mika}{οὔ τοι μὰ τὼ θεὼ σὺ καταπροίξει λέγουσα ταυτί,}{Non, par les deux déesses, tu ne l'auras pas dit gratuitement.}
\vers[n=567]{}{}{ἀλλ᾽ ἐκποκιῶ σου τὰς ποκάδας.}{Je t'arracherai la toison.}
\vers[suite]{ΜΝ.}{Mnésiloque}{οὐ δὴ μὰ Δία σύ γ᾽ ἅψει.}{Non, par Zeus, tu ne me toucheras pas.}
\vers[n=568]{ΜΙ.}{Mika}{καὶ μὴν ἰδού.}{Eh bien, tiens !}
\vers[suite]{ΜΝ.}{Mnésiloque}{καὶ μὴν ἰδού.}{Eh bien, tiens !}
\vers[suite]{ΜΙ.}{Mika}{λαβὲ θοἰμάτιον Φιλίστη.}{Prends mon manteau, Philista.}
\vers[n=569]{ΜΝ.}{Mnésiloque}{πρόσθιγε μόνον, κἀγώ σε νὴ τὴν Ἄρτεμιν—}{Atteins-moi seulement, et par Artémis...}
\vers[suite]{ΜΙ.}{Mika}{τί δράσεις;}{Que feras-tu ?}
\vers[n=570]{ΜΝ.}{Mnésiloque}{τὸν σησαμοῦνθ᾽ ὃν κατέφαγες, τοῦτον χεσεῖν ποιήσω.}{Le gâteau de sésame que tu as dévoré, je te le ferai évacuer.}
\vers[n=571]{ΚΟ.}{Coryphée}{παύσασθε λοιδορούμεναι· καὶ γὰρ γυνή τις ἡμῖν}{Trêve d'invectives. Car une femme}
\vers[n=572]{}{}{ἐσπουδακυῖα προστρέχει. πρὶν οὖν ὁμοῦ γενέσθαι,}{accourt vers nous en toute hâte. Avant qu'elle soit proche,}
\vers[n=573]{}{}{σιγᾶθ᾽, ἵν᾽ αὐτῆς κοσμίως πυθώμεθ᾽ ἅττα λέξει.}{faites silence. Écoutons en bonne contenance ce qu'elle va nous dire.}
\vers[n=574]{ΚΛ.}{Clisthène}{φίλαι γυναῖκες ξυγγενεῖς τοὐμοῦ τρόπου·}{Femmes aimées, qui m'êtes unies par la parenté des mœurs,}
\vers[n=575]{}{}{ὅτι μὲν φίλος εἴμ᾽ ὑμῖν, ἐπίδηλος ταῖς γνάθοις·}{que je vous suis ami, mes joues le prouvent assez.}
\vers[n=576]{}{}{γυναικομανῶ γὰρ προξενῶ θ᾽ ὑμῶν ἀεί.}{Oui, j'aime les femmes à la rage, et toujours je m'emploie pour vous.}
\vers[n=577]{}{}{καὶ νῦν ἀκούσας πρᾶγμα περὶ ὑμῶν μέγα}{Aussi, ayant appris sur votre sujet une chose capitale,}
\vers[n=578]{}{}{ὀλίγῳ τι πρότερον κατ᾽ ἀγορὰν λαλούμενον,}{dont on jasait tout à l'heure sur l'agora,}
\vers[n=579]{}{}{ἥκω φράσων τοῦτ᾽ ἀγγελῶν θ᾽ ὑμῖν, ἵνα}{je suis venu la raconter et vous donner avis}
\vers[n=580]{}{}{σκοπῆτε καὶ τηρῆτε μὴ καὶ προσπέσῃ}{de veiller et d'être en garde, de crainte}
\vers[n=581]{}{}{ὑμῖν ἀφάρκτοις πρᾶγμα δεινὸν καὶ μέγα.}{qu'une grave et fâcheuse affaire ne vous tombe sur les bras au dépourvu.}
\vers[n=582]{ΚΟ.}{Coryphée}{τί δ᾽ ἔστιν ὦ παῖ; παῖδα γάρ σ᾽ εἰκὸς καλεῖν,}{Qu'y a-t-il mon enfant ? Car il est naturel que je t'appelle enfant,}
\vers[n=583]{}{}{ἕως ἂν οὕτως τὰς γνάθους ψιλὰς ἔχῃς.}{tant que tu auras ainsi les joues lisses.}
\vers[n=584]{ΚΛ.}{Clisthène}{Εὐριπίδην φάσ᾽ ἄνδρα κηδεστήν τινα}{On dit qu'Euripide a envoyé ici même, aujourd'hui,}
\vers[n=585]{}{}{αὑτοῦ γέροντα δεῦρ᾽ ἀναπέμψαι τήμερον.}{un de ses parents, un vieillard.}
\vers[n=586]{ΚΟ.}{Coryphée}{πρὸς ποῖον ἔργον ἢ τίνος γνώμης χάριν;}{Pourquoi faire, dans quel dessein ?}
\vers[n=587]{ΚΛ.}{Clisthène}{ἵν᾽ ἅττα βουλεύοισθε καὶ μέλλοιτε δρᾶν,}{À cette fin que, quoi que vous agitiez ou décidiez de faire,}
\vers[n=588]{}{}{ἐκεῖνος εἴη τῶν λόγων κατάσκοπος.}{il puisse épier vos discours.}
\vers[n=589]{ΚΟ.}{Coryphée}{καὶ πῶς λέληθεν ἐν γυναιξὶν ὢν ἀνήρ;}{Et comment s'est-il dissimulé au milieu de femmes ? Un homme !}
